\documentclass[UTF8,fontset=fandol,11pt]{ctexart}
\usepackage[a4paper,margin=24mm]{geometry}
\usepackage{amsmath,amssymb,amsthm,graphicx,tabularx,booktabs,hyperref}
\newtheorem{definition}{定义}[section]
\newtheorem{theorem}[definition]{定理}
\newtheorem{example}[definition]{例题}
\newtheorem{exercise}[definition]{练习}
\hypersetup{unicode=true,colorlinks=true,linkcolor=blue}
\title{最小二乘法讲义计划}
\date{}
\begin{document}
\maketitle
\section{一元最小二乘拟合：从平方误差到正规方程}\label{sec:S01}
% UNIT S01.U01
\begin{definition}
设有 $n\geq 2$ 个观测点 $(x_i,y_i)$，其中 $x_i,y_i\in\mathbb{R}$，$i=1,\ldots,n$。用直线
\[
\hat y_i=a+bx_i
\]
拟合被解释变量 $y_i$，其中 $a,b\in\mathbb{R}$ 分别为截距参数和斜率参数。定义平方误差目标函数为
\begin{equation}
S(a,b)=\sum_{i=1}^{n}\bigl[y_i-(a+bx_i)\bigr]^2.
\label{def:ols-objective}
\end{equation}
一元最小二乘估计 $(\hat a,\hat b)$ 是使 $S(a,b)$ 在 $\mathbb{R}^2$ 上达到最小值的参数对。
\end{definition}
% UNIT S01.U02
\begin{definition}
对给定的参数 $a,b$，第 $i$ 个观测的残差定义为观测值与拟合值之差：
\begin{equation}
e_i=y_i-(a+bx_i).
\label{def:residual}
\end{equation}
当参数取最小二乘估计 $(\hat a,\hat b)$ 时，记
\[
\hat y_i=\hat a+\hat b x_i,
\qquad
\hat e_i=y_i-\hat y_i.
\]
此时平方误差目标函数的最小值为残差平方和 $\sum_{i=1}^n\hat e_i^2$。
\end{definition}
% UNIT S01.U03
\begin{theorem}
设 $x_i,y_i\in\mathbb{R}$，且
\label{thm:quadratic-unique-minimum}
\[
\sum_{i=1}^{n}(x_i-\bar x)^2>0,
\qquad
\bar x=\frac{1}{n}\sum_{i=1}^{n}x_i,
\qquad
\bar y=\frac{1}{n}\sum_{i=1}^{n}y_i.
\]
则平方误差函数 $S(a,b)$ 在 $\mathbb{R}^2$ 上具有唯一最小点 $(\hat a,\hat b)$。该最小点满足正规方程
\begin{equation}
\sum_{i=1}^{n}\bigl[y_i-(\hat a+\hat b x_i)\bigr]=0,
\qquad
\sum_{i=1}^{n}x_i\bigl[y_i-(\hat a+\hat b x_i)\bigr]=0,
\label{thm:normal-equations}
\end{equation}
并且
\[
\hat b=\frac{\sum_{i=1}^{n}(x_i-\bar x)(y_i-\bar y)}{\sum_{i=1}^{n}(x_i-\bar x)^2},
\qquad
\hat a=\bar y-\hat b\bar x.
\]
\end{theorem}
% UNIT S01.U04
\begin{proof}
将 $S(a,b)$ 展开得
\[
S(a,b)=\sum_{i=1}^{n}\bigl[y_i-a-bx_i\bigr]^2.
\]
对 $a$ 和 $b$ 分别求偏导，得到
\[
\frac{\partial S}{\partial a}=-2\sum_{i=1}^{n}\bigl[y_i-a-bx_i\bigr],
\qquad
\frac{\partial S}{\partial b}=-2\sum_{i=1}^{n}x_i\bigl[y_i-a-bx_i\bigr].
\]
因此，最小点必须满足式 \eqref{thm:normal-equations}。第一条正规方程给出
\[
n\hat a+\hat b\sum_{i=1}^{n}x_i=\sum_{i=1}^{n}y_i,
\]
从而
\[
\hat a=\bar y-\hat b\bar x.
\]
把它代入第二条正规方程，并利用
\[
\sum_{i=1}^{n}(x_i-\bar x)=0,
\qquad
\sum_{i=1}^{n}(y_i-\bar y)=0,
\]
可得
\[
\hat b\sum_{i=1}^{n}(x_i-\bar x)^2
=\sum_{i=1}^{n}(x_i-\bar x)(y_i-\bar y).
\]
由于分母严格为正，故 $\hat b$ 唯一确定，随后 $\hat a$ 也唯一确定。

下面验证由上述公式给出的 $(\hat a,\hat b)$ 确实满足正规方程。对第一条，有
\[
\sum_{i=1}^{n}\bigl[y_i-(\hat a+\hat b x_i)\bigr]
=\sum_{i=1}^{n}y_i-n\hat a-\hat b\sum_{i=1}^{n}x_i
=n\bar y-n(\bar y-\hat b\bar x)-\hat b n\bar x=0.
\]
对第二条，利用 $\sum_{i=1}^{n}(x_i-\bar x)=0$ 和 $\sum_{i=1}^{n}(y_i-\bar y)=0$，有
\[
\sum_{i=1}^{n}x_i\bigl[y_i-(\hat a+\hat b x_i)\bigr]
=\sum_{i=1}^{n}(x_i-\bar x)(y_i-\bar y)
-\hat b\sum_{i=1}^{n}(x_i-\bar x)^2=0,
\]
最后一步由 $\hat b$ 的定义得到。因此 $(\hat a,\hat b)$ 满足式 \eqref{thm:normal-equations}。

为说明这是唯一的全局最小点，令 $u=a-\hat a$，$v=b-\hat b$。由于 $(\hat a,\hat b)$ 满足正规方程，直接展开平方项可得
\[
S(a,b)=S(\hat a,\hat b)+\sum_{i=1}^{n}(u+vx_i)^2.
\]
再将 $x_i=(x_i-\bar x)+\bar x$ 代入，并令 $w=u+v\bar x$，得到
\[
\sum_{i=1}^{n}(u+vx_i)^2
=nw^2+v^2\sum_{i=1}^{n}(x_i-\bar x)^2\geq 0.
\]
当且仅当 $w=0$ 且 $v=0$ 时等号成立；由此 $u=0$。所以 $(\hat a,\hat b)$ 是唯一最小点。
\end{proof}
% UNIT S01.U05
\begin{theorem}
在定理 \eqref{thm:normal-equations} 的条件下，最小二乘拟合直线经过样本均值点 $(\bar x,\bar y)$，即
\[
\hat y(\bar x)=\hat a+\hat b\bar x=\bar y.
\]
等价地，含截距的一元最小二乘拟合满足
\[
\sum_{i=1}^{n}\hat e_i=0.
\]
\end{theorem}
% UNIT S01.U06
\begin{proof}
由定理 \eqref{thm:normal-equations} 给出的截距公式，
\[
\hat a=\bar y-\hat b\bar x.
\]
因此
\[
\hat a+\hat b\bar x=\bar y.
\]
另一方面，第一条正规方程为
\[
\sum_{i=1}^{n}\hat e_i=\sum_{i=1}^{n}\bigl[y_i-(\hat a+\hat b x_i)\bigr]=0\text{，}
\]
故拟合残差的总和为零。
\end{proof}
% UNIT S01.U07
\begin{example}
某地区三年的广告支出 $x$ 与销售额 $y$ 的观测数据为
\[
(x_1,y_1)=(1,2),\qquad (x_2,y_2)=(2,3),\qquad (x_3,y_3)=(3,5)\text{。}
\]
这三个观测点的样本均值为
\[
\bar x=\frac{1+2+3}{3}=2,\qquad \bar y=\frac{2+3+5}{3}=\frac{10}{3}\text{。}
\]
用一元直线拟合销售额与广告支出之间的线性关系。
\end{example}
% UNIT S01.U08
\par\noindent\textbf{解：}样本均值为
\[
\bar x=2,
\qquad
\bar y=\frac{10}{3}.
\]
因此
\[
\sum_{i=1}^{3}(x_i-\bar x)(y_i-\bar y)
=(-1)\left(-\frac43\right)+0\left(-\frac13\right)+1\left(\frac53\right)=3,
\]
且
\[
\sum_{i=1}^{3}(x_i-\bar x)^2=(-1)^2+0^2+1^2=2\text{。}
\]
所以
\[
\hat b=\frac{3}{2},
\qquad
\hat a=\frac{10}{3}-\frac32\cdot2=\frac13\text{。}
\]
拟合直线为
\[
\hat y=\frac13+\frac32x\text{。}
\]
三个观测点的拟合值和残差分别为
\[
\begin{array}{c|c|c}
 i&\hat y_i&\hat e_i=y_i-\hat y_i\\ \hline
 1&\frac{11}{6}&\frac16\\
 2&\frac{10}{3}&-\frac13\\
 3&\frac{29}{6}&\frac16
\end{array}
\]
可见 $\hat e_1+\hat e_2+\hat e_3=0$，且拟合直线经过 $(\bar x,\bar y)=(2,10/3)$。在经济解释上，该直线描述样本内广告支出与销售额的线性拟合关系，但不能仅凭拟合结果断言广告支出对销售额具有因果影响。\par
% UNIT S01.U09
\begin{exercise}
\label{mathagent:exercise:41706b9e371cbf8283a4f1192e660b4426644ffb8f650af52f47cef5491b1ca0}
设三家企业的培训支出 $x$ 与劳动生产率 $y$ 的数据为
\[
(x_1,y_1)=(1,4),\qquad (x_2,y_2)=(2,5),\qquad (x_3,y_3)=(4,8)\text{。}
\]
完成下列任务：
\begin{enumerate}
\item 计算 $\bar x$ 和 $\bar y$；
\item 求最小二乘拟合直线 $\hat y=\hat a+\hat b x$；
\item 计算三个观测点的残差，并验证残差之和为零；
\item 求培训支出为 $x=3$ 时的拟合值，并说明这是样本区间内的插值。
\end{enumerate}
\end{exercise}
\section*{练习答案}
% UNIT S01.U10
\par\noindent\textbf{\hyperref[mathagent:exercise:41706b9e371cbf8283a4f1192e660b4426644ffb8f650af52f47cef5491b1ca0]{练习 \ref*{mathagent:exercise:41706b9e371cbf8283a4f1192e660b4426644ffb8f650af52f47cef5491b1ca0} 的解答：}}答案如下：
\[
\bar x=\frac{1+2+4}{3}=\frac73,
\qquad
\bar y=\frac{4+5+8}{3}=\frac{17}{3}\text{。}
\]
进一步有
\[
\sum_{i=1}^{3}(x_i-\bar x)(y_i-\bar y)=\frac{19}{3},
\qquad
\sum_{i=1}^{3}(x_i-\bar x)^2=\frac{14}{3}\text{。}
\]
所以
\[
\hat b=\frac{19/3}{14/3}=\frac{19}{14},
\qquad
\hat a=\frac{17}{3}-\frac{19}{14}\cdot\frac73=\frac52\text{。}
\]
拟合直线为
\[
\hat y=\frac52+\frac{19}{14}x\text{。}
\]
三个拟合值和残差为
\[
\begin{array}{c|c|c}
 i&\hat y_i&\hat e_i=y_i-\hat y_i\\ \hline
 1&\frac{27}{7}&\frac17\\
 2&\frac{73}{14}&-\frac{3}{14}\\
 3&\frac{111}{14}&\frac{1}{14}
\end{array}
\]
残差之和为
\[
\frac17-\frac{3}{14}+\frac{1}{14}=0\text{。}
\]
最后，$x=3$ 时的拟合值为
\[
\hat y(3)=\frac52+\frac{19}{14}\cdot3=\frac{46}{7}\text{。}
\]
由于 $3$ 位于样本解释变量的最小值 $1$ 与最大值 $4$ 之间，这是样本区间内的插值。\par
\end{document}
